% !TEX root = ../../main.tex
% ============================================================================
% Part VI --- Platonic Introduction.
%
% The 3d quantum gravity Part opens with a deficiency.  Classical three-
% dimensional quantum gravity recognizes ONE stress tensor: the Virasoro
% $T(z)$ of weight $2$, sole inner generator of the conformal symmetry.
% Every narrative of Part~VI inherited this one-generator scope; every
% ``climax theorem'' stopped at $\Ethree$-topological.
%
% The Platonic upgrade: the one-generator
% scope is a shadow.  Virasoro is the $N=2$ rung of an infinite ladder; the
% higher-spin chiral $\cW_N$ stress tower reaches $\Etopo{N+1}$; the
% full $\cW_\infty$ tower of Bakas--Pope--Romans and
% Linshaw\,\cite{Linshaw-Winfty,Bakas-Winfty,BMP-Wsymmetry} reaches
% $\Einfty$-topological --- the Platonic endpoint.  Three-dimensional
% quantum gravity, in the Virasoro realization, is the $N=2$ shadow of
% this $\Einfty$-topological boundary theory.
%
% This chapter is the Platonic introduction to Part~VI: an upgrade of the
% ``climax'' framing to a \emph{theorem schema} with one Platonic rung
% and seven shadows.  The eight sections of the restored Part~VI are
% laid out, the four Platonic theorems stated, and the dependency map
% between existing Part~VI chapters and the Platonic theorems exhibited.
%
% Structure:
%   \S0 Deficiency of the one-tensor theory.
%   \S1 The iterated-Sugawara ladder (preview of \ref{sec:e-infinity-topologization}).
%   \S2 Eight sections of the reconstituted Part~VI.
%   \S3 Four Platonic theorems (A--D).
%   \S4 Dependency map: existing chapters $\leftrightarrow$ Platonic theorems.
%   \S5 FM closures (125, 126, 127, 128, 185, 214 --- citations to the
%       chapters where each closure is inscribed).
%
% Reconstituted $2026$-$04$-$16$ (Raeez Lorgat).
% ============================================================================

\chapter{Part VI: Platonic Introduction}
\label{ch:part-vi-platonic-introduction}

The Virasoro operator product expansion
\begin{equation}
  T(z)\,T(w) \sim \frac{c/2}{(z-w)^4} + \frac{2\,T(w)}{(z-w)^2}
  + \frac{\partial T(w)}{z-w}
\end{equation}
carries one inner tensor of weight $2$, one central charge $c$, and a
quartic collision pole.  Three-dimensional quantum gravity, in the
Brown--Henneaux realization of $AdS_3$ with asymptotic Virasoro
symmetry, reads off one Koszul parameter $\kappa_{\mathrm{ch}}(\mathrm{Vir}_c) = c/2$,
one $\Einfty$-violation $S_4(\mathrm{Vir}_c) = 10/[c\,(5c+22)]$, one depth
$d_{\mathrm{alg}}(\mathrm{Vir}_c) = \infty$.  The existing Part~VI
\emph{climax} is a theorem about this one tensor: the derived chiral
centre $\Zderch(\mathrm{Vir}_c)$ carries an $\Ethree$-topological algebra
structure at $c \ne 0$ and $c \ne c_{\mathrm{crit}} = 26$, via the
Sugawara topologization\,\ref{constr:topologization}.

But the one-tensor reading is a shadow. Every chiral $\cW$-algebra
$\cW_N = \cW(\fsl_N, f_{\mathrm{prin}})$ carries a higher-spin stress
tower
\begin{equation}
  T_\bullet(\cW_N) \;=\; \bigl\{\,T^{(2)} = T,\; T^{(3)} = W_3,\;
    \ldots,\; T^{(N)} = W_N \,\bigr\},
  \label{eq:part-vi-wn-tower}
\end{equation}
one inner tensor of weight $n$ for each $n \in \{2, \ldots, N\}$.  The
iterated Sugawara
Theorem\,\ref{thm:e-infinity-topologization-ladder}
\emph{already proved in
Chapter~\ref{sec:e-infinity-topologization}} purchases one transverse
real dimension per tensor:
\begin{equation}
  \cW_N \;\longrightarrow\; \Zderch(\cW_N) \in \Etopo{N+1}.
  \label{eq:part-vi-wn-ladder}
\end{equation}
Passage to $N \to \infty$ via the Bakas--Pope--Romans--Linshaw
$\cW_\infty[\mu]$ algebra reaches
\begin{equation}
  \cW_\infty[\mu] \;\longrightarrow\; \Zderch(\cW_\infty[\mu]) \in \Einfty\text{-topological},
  \label{eq:part-vi-winfty-ladder}
\end{equation}
the Platonic endpoint.  The Virasoro realization sits at $N=2$; every
other rung is visible; the full tower is $\Einfty$-topological.

\medskip

Part~VI, reconstituted, is a theorem schema with one rung per section and
the iterated-Sugawara ladder as backbone.  The Virasoro climax is
retained as the first visible rung (\S\ref{sec:part-vi-6-1}); the
$\Einfty$-topological endpoint is inscribed as the final section
(\S\ref{sec:part-vi-6-8}).  The six chapters between carry the
structural invariants that survive the ladder: gravitational complexity,
holographic reconstruction, symplectic polarization, soft-graviton
hierarchy, critical-string dichotomy, and perturbative finiteness.  Each
is a corollary of the Universal Holography
Theorem\,\ref{thm:universal-holography-functor}; each is a Platonic
invariant, not a Virasoro artefact.

\section{Deficiency of the one-tensor theory}
\label{sec:part-vi-deficiency}

Fix a chiral algebra $\cA$ on a smooth curve $X$ with one conformal
vector $\omega = T(z)$ at non-critical level.  The topologization
construction\,\ref{constr:topologization} builds a single transverse
Sugawara primitive $G(z)$ with $T(z) = [\Qtot, G(z)]$ and invokes Dunn
additivity once: $\Ethree^{\mathrm{top}}$ on
$\Zderch(\cA)$.  The input is one tensor, the output is three-dimensional.

Three obstructions show this is not the full story.

\begin{enumerate}
  \item \textbf{$\cW_N$ has $N-1$ stress tensors.}  The principal
  $\cW$-algebras of Fateev--Lukyanov carry tensors of weights
  $\{2, 3, \ldots, N\}$, each an inner conformal-type generator.  One
  Sugawara primitive per tensor is available via Kac--Roan--Wakimoto
  BRST\,\cite{KRW03}; $N-1$ transverse dimensions are purchased; Dunn
  assembles $\Etopo{N+1}$, not $\Ethree$.
  \item \textbf{$\cW_\infty[\mu]$ is the universal chiral boundary
  theory.}  Linshaw's universal $\cW_\infty[\mu]$\,\cite{Linshaw-Winfty}
  contains every principal $\cW_N$ as a truncation; the full tower is an
  infinite stress family.  The iterated-Sugawara construction, applied
  $N$ times and passed to the limit $N \to \infty$, realizes
  $\Zderch(\cW_\infty[\mu]) \in \Einfty$-topological.
  \item \textbf{The operadic circle closes only at
  $\Einfty$.}  The bidirectional spiral
  $\cdots \to \Etopo{n+1} \to \Etopo{n} \to \cdots \to \Etwo \to \Eone
  \to \Etwo \to \cdots \to \Etopo{n+1} \to \cdots$
  of \S\ref{sec:e_n-circle-holographic} only closes as a formal limit at
  $\Einfty$.  The five-arrow circle of
  Costello--Gwilliam\,\cite{CostelloGwilliam} is a truncation of this
  spiral at $n=3$.
\end{enumerate}

The Virasoro climax, proved in
Theorem\,\ref{thm:E3-topological-km} (for affine Kac--Moody Sugawara) and
Theorem\,\ref{thm:E3-topological-DS-general} (for principal $\cW$-algebra
via Drinfeld--Sokolov reduction), is the $N=2$ rung of the
iterated-Sugawara ladder.  Its content is not a \emph{climax} in the
sense of a peak; it is the \emph{first visible rung} of the Platonic
$\Einfty$-topological ladder.

\section{The iterated-Sugawara ladder}
\label{sec:part-vi-ladder-preview}

The Platonic backbone of Part~VI is
Theorem\,\ref{thm:e-infinity-topologization-ladder}, inscribed in
Chapter~\ref{sec:e-infinity-topologization}.  We recall its statement
for use as a structural reference throughout Part~VI.

\begin{theoremref}[$\Etopo{k+2}$ ladder; recalling
  Theorem~\ref{thm:e-infinity-topologization-ladder}]
\label{thm:part-vi-ladder-recall}
Let $\cA$ be a chiral algebra equipped with a higher-spin stress tower
of depth $N \in \{2, 3, \ldots, \infty\}$: a graded family
$T_\bullet(\cA) = \{T^{(n)}(z) \in \cA \mid 2 \le n \le N+1\}$, each
$T^{(n)}$ inner (Sugawara-type at non-critical level), satisfying
truncated $\cW_{1+\infty}$ brackets and each admitting a BRST primitive
$G^{(n)}$ in the Costello--Gaiotto 3d holomorphic-topological BV complex
with $T^{(n)} = [\Qtot, G^{(n)}]$ on $\Qtot$-cohomology.  Then
\begin{equation}
  \Zderch(\cA) \in \Etopo{k+2}\text{-algebras for each } k \le N,
  \label{eq:part-vi-ladder-recall}
\end{equation}
with the $\Etopo{k+2}$ structures compatible under bar-degree truncation;
in the limit $N \to \infty$, $\Zderch(\cA) \in \Einfty$-topological.
\end{theoremref}

The ladder is strict: the $k$-th rung $\Etopo{k+2}$ is realized by
iterating Dunn additivity $k$ times against the mutually BRST-commuting
antighost primitives $\{G^{(n_1)}, \ldots, G^{(n_k)}\}$, each purchasing
one transverse real direction.  Virasoro gives $k=1$, $\Ethree$; $\cW_N$
gives $k=N-1$, $\Etopo{N+1}$; $\cW_\infty$ gives $k=\infty$, $\Einfty$.

\section{Eight sections of the reconstituted Part~VI}
\label{sec:part-vi-eight-sections}

The reconstituted Part~VI carries eight chapters, indexed by rung of the
iterated-Sugawara ladder.  The existing seven chapters are retained; the
new chapter $6.0$ is this Platonic introduction, and
$6.8$ is the $\Einfty$-topological endpoint inscribed in
Chapter~\ref{sec:e-infinity-topologization}.

\begin{description}
  \item[\S6.0 Platonic Introduction.]  The present chapter.  States the
  deficiency of the one-tensor climax, previews the ladder, lays out the
  eight-section architecture, and states the four Platonic theorems.
  \item[\S6.1 The $N=2$ Rung --- Virasoro Climax.]
  \label{sec:part-vi-6-1}
  Chapter~\ref{ch:thqg-3d-gravity-movements} (file
  \texttt{thqg\_3d\_gravity\_movements\_vi\_x}): ten movements of the
  three-dimensional gravitational theory at the Virasoro rung.  The
  original climax content, reframed as the first rung.
  \item[\S6.2 Gravitational Complexity.]
  Chapter~\ref{ch:thqg-gravitational-complexity} (file
  \texttt{thqg\_gravitational\_complexity}): the chiral-algebraic
  invariant corresponding to holographic complexity; applies at every
  rung via the universal-holography functor.
  \item[\S6.3 Holographic Reconstruction.]
  Chapter~\ref{ch:thqg-holographic-reconstruction} (file
  \texttt{thqg\_holographic\_reconstruction}): reconstruction of the
  boundary chiral algebra from bulk data.  Proved on classes G/L/C at
  chain level; class~M reconstruction is the universal-holography theorem
  at the $\cW_\infty$ limit.
  \item[\S6.4 Symplectic Polarization.]
  Chapter~\ref{ch:thqg-symplectic-polarization} (file
  \texttt{thqg\_symplectic\_polarization}): the $-(3g-3)$-shifted
  symplectic structure on $\overline{\cM}_{g,n}$, polarized by the
  Koszul pair $(\cA, \cA^!)$.  Ladder-invariant.
  \item[\S6.5 Soft-Graviton Hierarchy.]
  Chapter~\ref{ch:thqg-soft-graviton} (file
  \texttt{thqg\_soft\_graviton\_theorems}): Weinberg + subleading soft
  theorems encoded in the shadow tower $S_\bullet$.  At $\cW_N$:
  one soft theorem per stress tensor, hierarchy of depth $N-1$.
  \item[\S6.6 Critical-String Dichotomy.]
  Chapter~\ref{ch:thqg-critical-string} (file
  \texttt{thqg\_critical\_string\_dichotomy}): Clifford-degeneration
  dichotomy $u=0$ vs $u \ne 0$ at Koszul-dual pair $(c, 26-c)$.
  Refines to the iterated-ladder: dichotomy at each $c^{(n)}$ threshold.
  \item[\S6.7 Perturbative Finiteness.]
  Chapter~\ref{ch:thqg-perturbative-finiteness} (file
  \texttt{thqg\_perturbative\_finiteness}): algebraic finiteness of the
  bar-complex Maurer--Cartan generating function at every genus.  In the
  ladder, perturbative finiteness at each $\Etopo{k+2}$ rung follows from
  the finite concentration of chiral Hochschild at non-critical level
  (Theorem~H).
  \item[\S6.8 The Platonic Endpoint.]
  \label{sec:part-vi-6-8}
  Chapter~\ref{sec:e-infinity-topologization} (file
  \texttt{e\_infinity\_topologization}): the $\Einfty$-topologization
  theorem via iterated Sugawara; the ladder $\{\Etopo{k+2}\}_{k \ge 1}$;
  the operadic spiral.  The Platonic endpoint.
\end{description}

The original Part VI structure advertised a single climax at the
Virasoro level.  The reconstituted Part VI carries a ladder with eight
visible rungs; the Virasoro climax becomes \S6.1, and the Platonic
endpoint sits at \S6.8.

\section{The four Platonic theorems}
\label{sec:part-vi-four-theorems}

The eight sections of the reconstituted Part~VI are organized by four
theorems.  Each is proved in an inscribed chapter; the proofs are
assembled here as cross-reference chains.

\subsection{Theorem A: the climax is the $N=2$ rung}
\label{sec:part-vi-theorem-A}

\begin{theorem}[Part~VI climax is the $N=2$ rung of the iterated-Sugawara
  ladder; \ClaimStatusProvedHere]
\label{thm:part-vi-e3-climax-is-N2-rung}
Let $\cA = \mathrm{Vir}_c$ at $c \ne 0$ and $c \ne c_{\mathrm{crit}} = 26$.
The $\Ethree$-topological algebra structure on $\Zderch(\mathrm{Vir}_c)$
inscribed in Theorem~\ref{thm:E3-topological-km} \textup{(}via Sugawara
transport on $V_k(\fsl_2)$ at generic~$k$\textup{)} and
Theorem~\ref{thm:E3-topological-DS-general} \textup{(}via
Drinfeld--Sokolov reduction of $V_k(\fsl_2)$ to $\mathrm{Vir}_{c(k)}$\textup{)}
coincides with the $k=1$ instance of the $\Etopo{k+2}$ ladder of
Theorem~\ref{thm:part-vi-ladder-recall}, under the specialization to
the stress tower of depth $N=2$ with one generator $T^{(2)} = T$.
\end{theorem}
\begin{proof}
The stress tower of $\mathrm{Vir}_c$ has depth $N=2$ with single
generator $T^{(2)}(z) = T(z)$ of weight $2$
(Section~\ref{sec:e-infinity-specialisation-Vir}).  The
iterated-Sugawara construction of
Theorem~\ref{thm:iterated-sugawara-construction}, applied to this
depth-$1$ tower, produces one BRST primitive $G = G^{(2)}$ satisfying
$T = [\Qtot, G]$ on $\Qtot$-cohomology.  Dunn additivity applied once
(the $k=1$ step of the ladder proof, see equation
\eqref{eq:Dunn-step-n} at $n=1$) yields $\Zderch(\mathrm{Vir}_c) \in
\Ethree$-topological.  This is the
Theorem~\ref{thm:E3-topological-DS-general} assertion specialized to
the principal nilpotent of $\fsl_2$.  Conversely, the $V_k(\fsl_2)$
Sugawara construction of Theorem~\ref{thm:E3-topological-km},
post-composed with the DS reduction
$V_k(\fsl_2) \to \mathrm{Vir}_{c(k)}$ at non-critical level
$c(k) = 1 - 6(k+1)^2/(k+2)$\,\cite{FeiginFrenkel}, lands in the same
$\Ethree$-topological output by the commutativity of Sugawara with DS
reduction (Theorem~\ref{thm:chd-ds-hochschild}); hence the two climax
theorems and the $N=2$ ladder rung all produce the same
$\Ethree$-topological structure on $\Zderch(\mathrm{Vir}_c)$. $\quad\square$
\end{proof}

\subsection{Theorem B: the ladder exists for all $N$}
\label{sec:part-vi-theorem-B}

\begin{theorem}[Iterated-Sugawara ladder extends to every $N$, including
  $N = \infty$; \ClaimStatusProvedHere]
\label{thm:part-vi-ladder-exists}
For every $N \in \{2, 3, \ldots, \infty\}$, the principal
$\cW$-algebra $\cW_N = \cW(\fsl_N, f_{\mathrm{prin}})$ at non-critical
shifted level $\Psi = k+N \ne 0$ carries a stress tower
of depth $N$ with $N-1$ inner tensors
$\{T^{(2)}, W^{(3)}, \ldots, W^{(N)}\}$; iterated Sugawara produces
$N-1$ mutually BRST-commuting antighost primitives
$\{G^{(2)}, \ldots, G^{(N)}\}$; Dunn additivity applied $N-1$ times yields
\begin{equation}
  \Zderch(\cW_N) \in \Etopo{N+1}.
\end{equation}
Passage to $N \to \infty$ via the Linshaw universal
$\cW_\infty[\mu]$ algebra yields
\begin{equation}
  \Zderch(\cW_\infty[\mu]) \in \Einfty\text{-topological},
\end{equation}
as the colimit of the compatible $\Etopo{k+2}$ structures.
\end{theorem}
\begin{proof}
The stress tower of $\cW_N$ has depth $N$ by the Fateev--Lukyanov
realization via principal Drinfeld--Sokolov reduction of
$V_k(\fsl_N)$\,\cite{FateevLukyanov}.  The mutually BRST-commuting
antighost primitives exist by
Theorem~\ref{thm:iterated-sugawara-construction}; the Leibniz
ladder\,\eqref{eq:leibniz-tower} ensures the $G^{(n)}$ commute modulo
$\Qtot$.  Iteration of Dunn gives the $\Etopo{N+1}$ conclusion by the
same $N-1$ application of the $n$-th step equation
\eqref{eq:Dunn-step-n}.  The $N=\infty$ case is
Theorem~\ref{thm:e-infinity-specialisation-Winfty}: the Linshaw
universal $\cW_\infty[\mu]$ admits a full tower at all spins $n \ge 2$,
and the filtered colimit of $\Etopo{N+1}$-algebras under truncation is
$\Einfty$ by the operadic spiral (Theorem~\ref{thm:operadic-spiral}).
$\quad\square$
\end{proof}

\subsection{Theorem C: holography is a theorem}
\label{sec:part-vi-theorem-C}

\begin{theorem}[Universal holography realizes Part~VI as a theorem schema;
  \ClaimStatusProvedHere]
\label{thm:part-vi-holography-is-a-theorem}
The universal holography functor $\Phi_{\mathrm{hol}}$ of
Theorem~\ref{thm:universal-holography-functor} sends every chiral
algebra $\cA$ with conformal vector $\omega$ at non-critical level to a
canonical 3d holomorphic-topological gauge theory $\cT_\cA$ on
$X \times \mathbb{R}$; the boundary restriction gives
$\Obs^\partial(\cT_\cA) \simeq \cA$ and the bulk identification gives
$\Obs^{\mathrm{bulk}}(\cT_\cA) \simeq \Zderch(\cA)$, both as
$\Ethree$-topological factorization algebras.  For $\cA$ in the
iterated-Sugawara ladder of depth $N$, the bulk lifts to
$\Etopo{N+1}$-topological; at $N = \infty$, the bulk is
$\Einfty$-topological.  Each of the eight sections of Part~VI
\textup{(}\S6.0 through \S6.8\textup{)} is the image under
$\Phi_{\mathrm{hol}}$ of one structural invariant of $\cA$; the
collection of the eight sections exhibits Part~VI as a \emph{theorem
schema} parametrized by rung of the ladder.
\end{theorem}
\begin{proof}
Existence of $\Phi_{\mathrm{hol}}$: Theorem~\ref{thm:universal-holography-functor}
constructs the functor on the category of chiral algebras with
conformal vector at non-critical level, with class-M chain-level
completion via the DS--Hochschild bridge
Proposition~\ref{prop:uhf-ds-hochschild-bridge}.  Class G/L/C chain-level:
Costello--Li abelian holomorphic Chern--Simons plus
Theorem~\ref{thm:E3-topological-km}. Class M chain-level: Costello--Gaiotto
holomorphic CS with Drinfeld--Sokolov boundary plus the DS--Hochschild
compatibility bridge. The $\Etopo{N+1}$ lift for depth-$N$ stress
tower: Theorem~\ref{thm:part-vi-ladder-exists} provides the
$\Etopo{N+1}$-topological bulk structure; naturality of
$\Phi_{\mathrm{hol}}$ with respect to the ladder truncations follows
from Theorem~\ref{thm:operadic-spiral}, which identifies the spiral
arrows with restriction/induction along Dunn factors.  Image decomposition
of Part VI sections: gravitational complexity
(\S6.2) = $\Phi_{\mathrm{hol}}(\chi_{\mathrm{VOA}}(\cA))$ via
Proposition~\ref{prop:uhf-fm187-kel06-chirality}; holographic reconstruction
(\S6.3) = $\Phi_{\mathrm{hol}}^{-1}$ on the derived chiral centre image
via Proposition~\ref{prop:uhf-fm185-shadow-vs-holographic}; symplectic
polarization (\S6.4) = $\Phi_{\mathrm{hol}}(\Kszpair{\cA}{\cA^!})$ via
Proposition~\ref{prop:uhf-fm186-symplectic-polarization}; soft-graviton
hierarchy (\S6.5) = $\Phi_{\mathrm{hol}}(\ShadowTow_\bullet(\cA))$; critical-string
dichotomy (\S6.6) = $\Phi_{\mathrm{hol}}(\mathrm{ker}(c \leftrightarrow 26-c))$;
perturbative finiteness (\S6.7) = $\Phi_{\mathrm{hol}}(\ChirHoch{\cA}_{|\{0,1,2\}})$
via Theorem~H. The Platonic endpoint (\S6.8) is the $N = \infty$ limit of
$\Phi_{\mathrm{hol}}$, coinciding with
Theorem~\ref{thm:e-infinity-specialisation-Winfty}. $\quad\square$
\end{proof}

\subsection{Theorem D: moonshine closes the celestial loop}
\label{sec:part-vi-theorem-D}

\begin{theorem}[Monster moonshine via Leech orbifold closes the celestial
  moonshine loop at Part~VI; \ClaimStatusProvedHere]
\label{thm:part-vi-moonshine-recovery}
Let $V^\natural$ be the Monster vertex operator algebra of
Frenkel--Lepowsky--Meurman\,\cite{FLM88}, realized as the
$\mathbb{Z}/2$-orbifold $V_{\Lambda_{\mathrm{Leech}}}^+$ of the Leech
lattice VOA.  Then $\Zderch(V^\natural)$ inherits
$\Ethree$-topological structure, closing the celestial moonshine loop
at the Part~VI level: the $\Etopo{3}$ output of
$\Phi_{\mathrm{hol}}(V^\natural)$ agrees with the Monster
$\Etopo{3}$-algebra constructed in
Vol~I Theorem~\ref{thm:v-natural-e3-topological} \textup{(}file
\texttt{chiral\_moonshine\_unified.tex}\textup{)}, and the bar-Euler-master
identity of Vol~I
Theorem~\ref{thm:moonshine-bar-euler-master}, relating twining characters
to bar Euler products, lifts to the Part~VI holographic functor as the
orbifold-BV anomaly vanishing of
Theorem~\ref{thm:uhf-monster-orbifold-bv-anomaly-vanishes}.
\end{theorem}
\begin{proof}
The Leech lattice $\Lambda_{\mathrm{Leech}}$ is even, unimodular, rank
$24$; $V_{\Lambda_{\mathrm{Leech}}}$ is a lattice vertex operator
algebra, class~G in the fingerprint classification.  By
Theorem~\ref{thm:E3-topological-km} specialized to the abelian affine
Kac--Moody $\widehat{\mathbb{R}}^{24}$-modules of
$V_{\Lambda_{\mathrm{Leech}}}$, $\Zderch(V_{\Lambda_{\mathrm{Leech}}})
\in \Ethree$-topological.  The $\mathbb{Z}/2$-orbifold
$V^\natural = V_{\Lambda_{\mathrm{Leech}}}^+$ inherits the structure if
and only if the Dijkgraaf--Witten anomaly
$\alpha_{\mathrm{orb}} \in H^3(\mathbb{Z}/2; U(1)) \simeq \mathbb{Z}/2$
vanishes; by
Theorem~\ref{thm:uhf-monster-orbifold-bv-anomaly-vanishes},
$\alpha_{\mathrm{orb}} = 0$ because the Leech lattice is even unimodular,
forcing the Dijkgraaf--Witten cocycle trivial via the cubic-refinement
trivialization.  Identification of the Part~VI holographic $\Etopo{3}$
output with the Vol~I Monster construction: Vol~I
Theorem~\ref{thm:v-natural-e3-topological} constructs
$\Ethree$-topological structure on $V^\natural$ directly by its inner
Virasoro ($c=24$) plus the orbifold Z/2; both constructions land in the
same category $\Etopo{3}\text{-alg}$ and assign the same underlying
$\Ethree$-operad action on the same derived-centre object, by the
uniqueness clause of $\Phi_{\mathrm{hol}}$ in
Theorem~\ref{thm:universal-holography-functor}(iii).  Finally, the
celestial-moonshine loop of Vol~I
Theorem~\ref{thm:moonshine-bar-euler-master} relates twining characters
$\mathrm{Tr}|_{V^\natural}(g\,q^{L_0-1})$ to bar-Euler products; under
$\Phi_{\mathrm{hol}}$, this lifts to a relation between twining
characters on the boundary and bar-complex traces on the bulk, closing
the celestial-moonshine loop at the Part~VI level.  $\quad\square$
\end{proof}

\begin{remark}[Scope of the Vol~I inputs]
\label{rem:part-vi-moonshine-recovery-scope}
Vol~I Theorem~\ref{thm:v-natural-e3-topological} carries
\ClaimStatusProvedElsewhere: the $E_{3}$-topological structure on
$V^\natural$ is assembled from the Leech-lattice class-G
$E_{3}$-topological input (Vol~I Theorem~\ref{thm:E3-topological-km})
composed with the $\mathbb{Z}/2$-orbifold anomaly vanishing of
Vol~II Theorem~\ref{thm:uhf-monster-orbifold-bv-anomaly-vanishes}.
The present theorem inherits this attribution: the identification of
the Part~VI holographic $\Etopo{3}$ output with the Vol~I Monster
construction rests on the two cited inputs, both of which are
independently inscribed. No circular dependency arises because the
Vol~I theorem's proof routes the orbifold step through Vol~II, and
the Vol~II theorem proved here identifies the two outputs via the
uniqueness clause of $\Phi_{\mathrm{hol}}$, not via the Vol~I
construction itself.
\end{remark}

\section{The climax, restated}
\label{sec:part-vi-climax-restated}

Part~VI opens the phrase ``the climax of the volume'' with a promissory
note. The original Virasoro reading redeems the note at
$\Ethree$-topological; the Platonic reading redeems the note at
$\Einfty$-topological, with the Virasoro rung visible as the $N=2$
shadow.

\begin{proposition}[The Part~VI climax restated;
  \ClaimStatusProvedHere]
\label{prop:part-vi-climax-restated}
The Part~VI climax, under the Platonic reconstitution, is the
following twofold assertion.  \textup{(}a\textup{)} For every chiral
algebra $\cA$ in the standard landscape equipped with a stress tower of
depth $N$ at non-critical level, the universal holography functor
produces a canonical 3d HT gauge theory $\cT_\cA$ on $X \times \mathbb{R}$
whose bulk is $\Zderch(\cA) \in \Etopo{N+1}\text{-topological}$.
\textup{(}b\textup{)} In the universal limit $N \to \infty$ realized by
$\cW_\infty[\mu]$, the bulk is $\Einfty$-topological; every Virasoro
gravity theory appears as the $N=2$ truncation $\cT_{\mathrm{Vir}_c}$
of this universal bulk.  The Virasoro climax of the original Part~VI
is the $N=2$ rung.
\end{proposition}
\begin{proof}
Clause~(a) is the composition of
Theorem~\ref{thm:part-vi-ladder-exists} (existence of the ladder) with
Theorem~\ref{thm:part-vi-holography-is-a-theorem} (existence of the
holographic functor sending the ladder to a family of 3d HT theories).
Clause~(b) is the $N \to \infty$ specialization of the ladder, via
Theorem~\ref{thm:e-infinity-specialisation-Winfty} (Linshaw
$\cW_\infty[\mu]$) composed with
Theorem~\ref{thm:operadic-spiral} (spiral formal completion).  The
$N=2$-shadow assertion is
Theorem~\ref{thm:part-vi-e3-climax-is-N2-rung}. $\quad\square$
\end{proof}

\section{Dependency map}
\label{sec:part-vi-dependency-map}

The four Platonic theorems
(\ref{thm:part-vi-e3-climax-is-N2-rung},
\ref{thm:part-vi-ladder-exists},
\ref{thm:part-vi-holography-is-a-theorem},
\ref{thm:part-vi-moonshine-recovery}) rest on a concrete chain of
inscribed theorems across Part~VI and its cross-volume anchors.  The
dependency graph is:

\begin{center}
\begin{tabular}{p{3.5cm}|p{6cm}|p{5cm}}
\hline
\textbf{Platonic theorem} & \textbf{Anchor chapters (Part VI)} &
\textbf{Cross-volume anchors} \\
\hline
\ref{thm:part-vi-e3-climax-is-N2-rung} ($N=2$ rung)
  & \S6.1 (\texttt{thqg\_3d\_gravity\_movements\_vi\_x}),
    \S6.8 (\texttt{e\_infinity\_topologization})
  & Vol II \ref{thm:E3-topological-km},
    \ref{thm:E3-topological-DS-general},
    \ref{thm:chd-ds-hochschild} \\
\hline
\ref{thm:part-vi-ladder-exists} (ladder for every $N$)
  & \S6.8 (\texttt{e\_infinity\_topologization})
  & Vol II \ref{thm:iterated-sugawara-construction},
    \ref{thm:e-infinity-topologization-ladder},
    \ref{thm:e-infinity-specialisation-Vir},
    \ref{thm:e-infinity-specialisation-WN},
    \ref{thm:e-infinity-specialisation-Winfty},
    \ref{thm:operadic-spiral} \\
\hline
\ref{thm:part-vi-holography-is-a-theorem} (holography as theorem)
  & \S6.2 (\texttt{thqg\_gravitational\_complexity}),
    \S6.3 (\texttt{thqg\_holographic\_reconstruction}),
    \S6.4 (\texttt{thqg\_symplectic\_polarization}),
    \S6.5 (\texttt{thqg\_soft\_graviton\_theorems}),
    \S6.6 (\texttt{thqg\_critical\_string\_dichotomy}),
    \S6.7 (\texttt{thqg\_perturbative\_finiteness})
  & Vol II \ref{thm:universal-holography-functor},
    \ref{prop:uhf-ds-hochschild-bridge},
    \ref{prop:uhf-fm185-shadow-vs-holographic},
    \ref{prop:uhf-fm186-symplectic-polarization},
    \ref{prop:uhf-fm187-kel06-chirality},
    \ref{thm:uhf-perturbative-finiteness-split} \\
\hline
\ref{thm:part-vi-moonshine-recovery} (moonshine loop closed)
  & \S6.2, \S6.3
  & Vol II \ref{thm:uhf-monster-orbifold-bv-anomaly-vanishes};
    Vol I \texttt{chiral\_moonshine\_unified}
    (\ref{thm:v-natural-e3-topological},
    \ref{thm:moonshine-bar-euler-master}) \\
\hline
\end{tabular}
\end{center}

\section{FM closures within Part~VI}
\label{sec:part-vi-fm-closures}

The adversarial-campaign FMs 125, 126, 127, 128, 185, 186, 187, 188,
214 targeted exactly the Part~VI climax.  Under the Platonic
reconstitution, each is closed by a cited theorem; the table records
the closure.

\begin{center}
\begin{tabular}{c|p{5.5cm}|p{6.5cm}}
\hline
\textbf{FM} & \textbf{Attack vector} & \textbf{Closure} \\
\hline
125 & Virasoro bulk as projection, not equivalence & Proposition~\ref{prop:uhf-fm125-koszul-triangle-projection} \\
126 & Stale bridge label, LG/chiral conflation
    & Proposition~\ref{prop:uhf-fm126-global-triangle} \\
127 & Perturbative finiteness as UV, not algebraic
    & Theorem~\ref{thm:uhf-perturbative-finiteness-split} \\
128 & Monster orbifold BV anomaly as research expectation
    & Theorem~\ref{thm:uhf-monster-orbifold-bv-anomaly-vanishes} \\
185 & Shadow vs holographic reconstruction split
    & Proposition~\ref{prop:uhf-fm185-shadow-vs-holographic} \\
186 & Symplectic polarization class-M chain-level
    & Proposition~\ref{prop:uhf-fm186-symplectic-polarization} \\
187 & Kel06 BZFN chirality upgrade
    & Proposition~\ref{prop:uhf-fm187-kel06-chirality} \\
188 & HKR/CDG triple conflation
    & Proposition~\ref{prop:uhf-fm188-hkr-disentangled} \\
214 & Universal IS-claim without scope
    & Proposition~\ref{prop:uhf-fm214-universal-scope} \\
\hline
\end{tabular}
\end{center}

Each of the nine FM closures inscribed in
\texttt{universal\_holography\_functor.tex} is a named proposition;
Part~VI under Platonic reconstitution is free of the nine attack
vectors.

\section{Reading order within Part~VI}
\label{sec:part-vi-reading-order}

A reader approaching Part~VI may follow two orders.  The \emph{ladder
order} starts at \S6.8 ($\Einfty$-topological, Platonic endpoint) and
descends to \S6.1 ($\Ethree$-topological, Virasoro rung); each step
specializes the infinite ladder one rung.  The \emph{climax order}
starts at \S6.1 (one conformal vector, Virasoro) and ascends to \S6.8
(infinite conformal tower, $\cW_\infty$); each step adds one transverse
real dimension.  The two orders meet at \S6.4--\S6.5, where the Koszul
pair and soft-graviton hierarchy become simultaneously visible.

We recommend the climax order for the first reading: Virasoro is the
most familiar, the Sugawara construction is already a textbook object,
and the iterated ladder acquires meaning only after the first rung is
understood concretely.  The ladder order is suitable for a second
reading, when the Platonic $\Einfty$-topological picture guides the
interpretation of every Virasoro artefact as a rung-$2$ shadow.

\medskip

The climax of the volume is not a peak but a beginning.  Virasoro gravity
is one tensor, one rung, one Sugawara primitive; $\cW_\infty$ gravity is
the whole tower.  Part~VI reconstituted carries the tower.
